% !TEX root = ../main.tex
\chapter{紧致性、连续性与 Weierstrass 极值定理}
\label{ch:compact}

\noindent\emph{用到的前置工具：度量空间、开集与闭集、序列收敛与完备性（第八章）；上确界与下确界、实数的完备性（实数公理，线性代数与分析部分常用）。}
\medskip

经济学里几乎每一个模型，第一步都是写下一个最优化问题：消费者最大化效用、厂商最小化成本、计量学家在参数空间上最大化对数似然。可在动笔求一阶条件之前，有一个更基本、却常被跳过的问题：\textbf{这个最优解到底存不存在？}如果连``最大值''都不存在，那么``最优选择''、``均衡''、``估计量''就都是空中楼阁，后面所有比较静态、所有渐近分析都失去了立足点。本章给出回答这个问题的标准工具——\textbf{Weierstrass 极值定理}：连续函数在紧致集上一定取到最大值与最小值。为把它讲清楚，我们要先认识两个角色——刻画定义域``不漏气''的\textbf{紧致性}，和刻画目标函数``不跳变''的\textbf{连续性}——再看它们如何合力保证最优解存在。最后简单提一句一致连续，它是后续讲积分与逼近时会再登场的近亲。

\section{先问``解存不存在''}
\label{sec:compact-motivation}

考虑最朴素的最大化问题
\[
    \max_{x\in S} f(x),
\]
$S$ 是可行集（预算集、参数空间……），$f$ 是目标函数（效用、利润、似然）。我们太习惯``求出最优解''这件事，以至于默认它一定有解。但它\emph{未必}有。

举两个一眼可见的失败。其一，把可行集取成开区间 $S=(0,1)$、目标 $f(x)=x$：$f$ 在 $S$ 上的取值可以任意接近 $1$，却永远到不了 $1$——没有哪个 $x\in(0,1)$ 使 $f(x)$ 最大，``最优解''不存在。毛病出在定义域``漏了气'':它缺了本该是上确界所在的那个端点。其二，把可行集取成整条实轴 $S=\RR$、目标 $f(x)=x$：取值无上界，更谈不上最大。毛病出在定义域``跑到了无穷远''。这两种病——\emph{缺边界}与\emph{跑无穷}——正是``紧致''这个概念要一举堵死的。

还有第三种失败，毛病不在定义域而在函数自身：哪怕 $S=[0,1]$ 这么规整，若 $f$ 在某点突然跳一下（不连续），它也可能在那一点附近无限接近一个值却取不到。把这三种失败合起来看，结论很清爽：\emph{只要定义域紧致、函数连续，三种病就全被排除，最大值与最小值必定存在}。这就是 Weierstrass 极值定理，本章的主角。

\state{为什么经济学家必须在乎``存在性''}{
``最优解存在''不是数学洁癖，而是整座建筑的地基。需求函数 $x^*(p,w)=\argmax_{x\in B(p,w)}u(x)$ 要先有解，才谈得上它怎么随价格变（比较静态）；一般均衡要先有解，才谈得上它的性质；极大似然估计量 $\hat\theta=\argmax_\theta \ell_n(\theta)$ 要先\emph{良定义}，才谈得上它一致、渐近正态。本章的定理，就是这一切``先有解''的统一出处。
}

\section{紧致性}
\label{sec:compact-compactness}

紧致性有两副面孔：一副是抽象的``开覆盖''定义，它在任意拓扑空间里都成立、是定理证明里最锋利的工具；另一副是直观的``有界闭集''刻画，它只在 $\RR^n$ 里成立、却是我们日常判断一个集合紧不紧致的实用判据。两者之间架着 Heine--Borel 定理这座桥。

\subsection{开覆盖定义}

\defn{开覆盖与紧致集}{
设 $K$ 是度量空间 $(X,d)$ 的子集。$X$ 中一族开集 $\set{U_\alpha}_{\alpha\in A}$ 称为 $K$ 的一个\textbf{开覆盖}，若
\[
    K\;\se\;\bigcup_{\alpha\in A} U_\alpha .
\]
若从中总能挑出\emph{有限}个 $U_{\alpha_1},\dots,U_{\alpha_n}$ 仍覆盖 $K$（即 $K\se U_{\alpha_1}\cup\cdots\cup U_{\alpha_n}$），就称这是一个\textbf{有限子覆盖}。若 $K$ 的\emph{每一个}开覆盖都存在有限子覆盖，则称 $K$ 是\textbf{紧致集}（compact）。
}

这个定义初看抽象得让人摸不着头脑：为什么``任意开覆盖都能缩成有限个''会是个有用的性质？关键在那个``任意''与``有限''的对比。它说的是：无论别人用多么稀碎、多么刁钻的一族开集来盖住 $K$，你总能只留下有限件就把活干完（图~\ref{fig:compact-1}）。这种``无穷的覆盖可以被有限地驾驭''的能力，正是后面把``逐点的局部信息''拼成``整体结论''时所依赖的——证明 Weierstrass 定理、证明连续函数在紧集上一致连续，用的都是它。

\begin{figure}[h]
\centering
\begin{tikzpicture}[xscale=1.5,yscale=1.15,>=Stealth]
  % 数轴
  \draw[->] (-0.3,0) -- (5.6,0) node[right] {$\RR$};
  % 闭区间 [a,b]=[0.5,4.5]，加粗，端点实心
  \draw[very thick,blue!60!black] (0.5,0) -- (4.5,0);
  \fill[blue!60!black] (0.5,0) circle (1.6pt);
  \fill[blue!60!black] (4.5,0) circle (1.6pt);
  \node[blue!60!black,anchor=north] at (0.5,-0.07) {$a$};
  \node[blue!60!black,anchor=north] at (4.5,-0.07) {$b$};
  \node[blue!60!black,anchor=north] at (2.5,-0.34) {$K=[a,b]$};
  % 一族开区间（覆盖），两个高度交错，端点小竖杠表示开端
  \draw[red!70!black,thick] (0.25,0.50) -- (1.65,0.50);
  \draw[red!70!black,thick] (0.25,0.45) -- (0.25,0.55);
  \draw[red!70!black,thick] (1.65,0.45) -- (1.65,0.55);
  \node[red!70!black,anchor=south,font=\small] at (0.95,0.53) {$U_1$};
  \draw[red!70!black,thick] (2.40,0.50) -- (3.85,0.50);
  \draw[red!70!black,thick] (2.40,0.45) -- (2.40,0.55);
  \draw[red!70!black,thick] (3.85,0.45) -- (3.85,0.55);
  \node[red!70!black,anchor=south,font=\small] at (3.12,0.53) {$U_3$};
  \draw[red!70!black,thick] (1.30,0.95) -- (2.75,0.95);
  \draw[red!70!black,thick] (1.30,0.90) -- (1.30,1.00);
  \draw[red!70!black,thick] (2.75,0.90) -- (2.75,1.00);
  \node[red!70!black,anchor=south,font=\small] at (2.02,0.98) {$U_2$};
  \draw[red!70!black,thick] (3.50,0.95) -- (4.90,0.95);
  \draw[red!70!black,thick] (3.50,0.90) -- (3.50,1.00);
  \draw[red!70!black,thick] (4.90,0.90) -- (4.90,1.00);
  \node[red!70!black,anchor=south,font=\small] at (4.20,0.98) {$U_4$};
  % 竖直细虚线把开区间端点投到数轴，示意覆盖
  \foreach \x in {0.25,1.65,2.40,3.85}{
     \draw[gray!40,dotted] (\x,0.04) -- (\x,0.44);
  }
  \foreach \x in {1.30,2.75,3.50,4.90}{
     \draw[gray!40,dotted] (\x,0.04) -- (\x,0.89);
  }
  % 说明文字：放最上方，留足净空
  \node[anchor=west,align=left,font=\small] at (0.05,1.78)
     {紧致 $=$ 任一开覆盖 $\{U_\alpha\}$ 都有有限子覆盖；};
  \node[anchor=west,align=left,font=\small] at (0.05,1.56)
     {此处取 $U_1,U_2,U_3,U_4$ 这四个就盖住了 $K$};
\end{tikzpicture}
\caption{开覆盖与有限子覆盖。一堆开区间盖住闭区间 $K=[a,b]$；紧致性断言，无论原来这族开集有多少个，总能只留下有限个（这里四个）仍盖住 $K$。}
\label{fig:compact-1}
\end{figure}

\rmkb[为什么用``开''集、为什么要``任意''覆盖]{
两处都不可放松。若允许用闭集覆盖，则任何集合都能被它自身（一个闭集）``覆盖''，定义就废了——开集的``软边界''才让覆盖这件事有内容。而``任意''也省不得：开区间 $(0,1)$ \emph{存在}有限子覆盖的开覆盖（例如用单个 $(-1,2)$ 盖住），但它\emph{不是}对每个开覆盖都行——取 $U_n=(1/n,\,1)$，$n=1,2,\dots$，这族开集盖住了 $(0,1)$，可任意有限个的并都是某个 $(1/N,1)$，漏掉了 $(0,1/N]$。正是这个``漏''，暴露出 $(0,1)$ 不紧致——也正是上一节``取不到最大值''那个毛病的拓扑版本。
}

\subsection{Heine--Borel：在 \texorpdfstring{$\RR^n$}{Rn} 里紧致就是有界闭}

开覆盖定义虽利，却难直接验证——你总不能真去检查``每一个''开覆盖。所幸在欧氏空间里，紧致性等价于一个一眼可查的几何条件。

\thm{Heine--Borel 定理}{
设 $K\se\RR^n$。则
\[
    K \text{ 紧致}
    \quad\Iff\quad
    K \text{ 既有界又闭}.
\]
}

\rmkb[这是 $\RR^n$ 的特权，不是普遍真理]{
``紧致 $=$ 有界闭''只在有限维欧氏空间（更一般地，在有限维赋范空间）里成立。在无穷维空间里它会崩塌：闭单位球在任何无穷维赋范空间里都是有界闭的，却\emph{不}紧致。这件事对计量与决策论并非纯属遐想——当参数空间是函数空间（非参数估计、连续型机制设计）时，``有界闭''就不再保证存在性，得另请高明（如 Arzelà--Ascoli 给出函数族紧致的判据）。本书后续凡说``紧致''，默认在 $\RR^n$ 中，可与``有界闭''互换。
}

完整证明偏重纯拓扑，这里只给两个方向的直觉，足够使用与记忆。

\pf[直觉，非严格]{
\textbf{紧致 $\To$ 有界闭。}先说有界：用一族同心开球 $B(0,m)$（$m=1,2,\dots$）覆盖 $K$，它们是一个开覆盖；有限子覆盖意味着 $K$ 被某个最大的 $B(0,M)$ 装下，故有界。再说闭：若 $K$ 不闭，则存在一个在 $K$ 外的极限点 $p$，可用``挖掉 $p$ 的小邻域''的那些开集覆盖 $K$；但因 $p$ 是极限点，任何有限个这样的开集都会在 $p$ 附近留下 $K$ 的点没盖住——与紧致矛盾。故 $K$ 必闭。

\smallskip
\noindent\textbf{有界闭 $\To$ 紧致。}这一半是实数完备性的化身。先看一维闭区间 $[a,b]$：用\emph{二分法}。若某个开覆盖没有有限子覆盖，就把 $[a,b]$ 平分，两半中至少一半仍不能被有限子覆盖；对它再平分，如此下去得到一串长度趋于零、层层套住的闭区间，每个都不能被有限子覆盖。由闭区间套定理（实数完备性），它们交于唯一一点 $x^*\in[a,b]$。但 $x^*$ 总被覆盖里某个开集 $U$ 含住，而 $U$ 是开的、必含 $x^*$ 的一个小邻域——当区间套缩得足够小时，整段区间都落进这一个 $U$ 里，它\emph{单独一个}就构成有限子覆盖，矛盾。故 $[a,b]$ 紧致。高维的有界闭集落在某个闭方体 $[a,b]^n$ 内，方体由一维情形（取乘积）紧致，而紧致集的闭子集仍紧致，于是 $K$ 紧致。
}

\state{$\RR^n$ 中判紧致的实用流程}{
要断定 $K\se\RR^n$ 紧致，只需查两件事：\textbf{(1) 有界}——存在 $R$ 使 $\norm{x}\le R$ 对一切 $x\in K$ 成立；\textbf{(2) 闭}——$K$ 的每个收敛序列的极限仍在 $K$ 内（等价地，$K$ 含其全部极限点）。预算集 $\setb{x\ge 0}{p\cdot x\le w}$（价格为正时）、概率单纯形 $\setb{x\ge 0}{\sum_i x_i=1}$、有界闭的参数空间 $\Theta$，都是这样验出来的紧致集——它们正是消费者问题、混合策略、极值估计存在性的舞台。
}

\subsection{序列紧致：紧致性的``可用'' 形态}

证明 Weierstrass 定理时，我们其实不直接搬开覆盖，而用紧致性的一个等价说法，它更贴合``序列''这门第八章已练熟的语言。

\thm{序列紧致（$\RR^n$ 中等价于紧致）}{
$K\se\RR^n$ 紧致，当且仅当 $K$ 是\textbf{序列紧致}的：$K$ 中\emph{任意}序列 $\set{x_k}$ 都有一个收敛子列 $\set{x_{k_j}}$，且其极限仍落在 $K$ 内。
}

\rmkb[为什么序列紧致就是``有界闭''的序列版]{
拆开看：``$K$ 中任意序列有收敛子列''正是 \textbf{Bolzano--Weierstrass 定理}（$\RR^n$ 中任何有界序列必有收敛子列）所保证的——前提是序列有界，而这由 $K$ 有界供给；``极限仍在 $K$ 内''则正是 $K$ \emph{闭}的定义。所以``有界''管``抽得出收敛子列''，``闭''管``极限不逃出 $K$''。两件事合起来，就是序列紧致。这也再次印证了 Heine--Borel：有界闭 $\Iff$ 序列紧致 $\Iff$ 紧致。
}

\section{连续性}
\label{sec:compact-continuity}

定义域备好了，再看目标函数。``连续''要捕捉的，是``输入微调、输出也只微调''这件直觉上``函数不跳变''的事。它同样有两副等价的面孔——古典的 $\ve$--$\delta$ 语言，和更便于和紧致性配合的序列语言。

\defn{连续（$\ve$--$\delta$ 与序列两种刻画）}{
设 $f:S\to\RR$，$S\se\RR^n$，$x_0\in S$。以下两条等价，满足任一条即称 $f$ 在 $x_0$ \textbf{连续}；若 $f$ 在 $S$ 的每一点都连续，称 $f$ 在 $S$ 上连续。
\begin{itemize}
    \item[(1)] \textbf{$\ve$--$\delta$ 刻画.} 对任意 $\ve>0$，存在 $\delta>0$，使得对一切 $x\in S$，
    \[
        \norm{x-x_0}<\delta
        \;\To\;
        \abs{f(x)-f(x_0)}<\ve .
    \]
    \item[(2)] \textbf{序列刻画.} 对 $S$ 中任意收敛到 $x_0$ 的序列 $\set{x_k}$，都有 $f(x_k)\to f(x_0)$。
\end{itemize}
}

\rmkb[两种刻画为何等价、各擅何处]{
$\ve$--$\delta$ 量化了``要输出误差小于 $\ve$，输入只需贴近到 $\delta$ 之内''，适合做精细的估计；序列刻画则把连续等同于``可以与取极限交换次序''——$\lim_k f(x_k)=f(\lim_k x_k)$，做收敛论证时格外顺手。证明二者等价的关键一步是：若序列刻画成立而 $\ve$--$\delta$ 不成立，便能造出一个收敛到 $x_0$、函数值却始终离 $f(x_0)$ 至少 $\ve$ 的序列，与序列刻画矛盾（这里悄悄用到了选取——对每个 $\delta=1/k$ 挑一个``坏点'' $x_k$）。下一节我们用的是序列刻画，因为它能和序列紧致严丝合缝地咬在一起。
}

\section{Weierstrass 极值定理}
\label{sec:compact-weierstrass}

万事俱备。把``定义域紧致''与``函数连续''合在一处，最大值与最小值的存在性就水到渠成。

\thm{Weierstrass 极值定理}{
设 $K\se\RR^n$ 非空紧致，$f:K\to\RR$ 连续。则 $f$ 在 $K$ 上\textbf{有界}，并且\textbf{达到}它的最大值与最小值：存在 $x^*,x_*\in K$ 使得
\[
    f(x_*)\;\le\;f(x)\;\le\;f(x^*)
    \qquad\text{对一切 } x\in K .
\]
等价地，$\max_{x\in K} f(x)$ 与 $\min_{x\in K} f(x)$ 都存在（而非仅有 $\sup$、$\inf$）。
}

证明分两步，逻辑很干净：先证``有界''（于是上确界是个有限数），再证这个上确界``真被取到''。两步都靠序列紧致把``逼近''升级成``达到''。最小值的论证与最大值对称，证一个即可。

\pf{
\textbf{第一步：$f$ 有上界。}反证。若 $f$ 在 $K$ 上无上界，则对每个 $k\in\NN$ 都能找到 $x_k\in K$ 使 $f(x_k)>k$。由 $K$ 序列紧致，$\set{x_k}$ 有收敛子列 $x_{k_j}\to \bar x\in K$。$f$ 连续，故由序列刻画 $f(x_{k_j})\to f(\bar x)$，是个\emph{有限}数。但按构造 $f(x_{k_j})>k_j\to\infty$，发散——矛盾。故 $f$ 有上界（同理有下界），$f$ 在 $K$ 上有界。

\smallskip
\noindent\textbf{第二步：上确界被取到。}既然 $f$ 在 $K$ 上有上界，记
\[
    M=\sup_{x\in K} f(x)\in\RR .
\]
由上确界的定义，存在一列 $x_k\in K$ 使 $f(x_k)\to M$（这样的``逼近最大值''的序列总能取到——这是 $\sup$ 的特征）。再用序列紧致：$\set{x_k}$ 有收敛子列 $x_{k_j}\to x^*\in K$。一方面，作为收敛子列，它对应的函数值仍 $f(x_{k_j})\to M$；另一方面，$f$ 连续给出 $f(x_{k_j})\to f(x^*)$。极限唯一，于是
\[
    f(x^*)=M=\sup_{x\in K} f(x) .
\]
这说明上确界 $M$ 在点 $x^*\in K$ 处\emph{真正被取到}，它就是最大值。对 $-f$ 重复整套论证（$-f$ 也连续），即得最小值在某 $x_*\in K$ 取到。
}

\state{定理的``三步骨架''值得记住}{
整个证明就是同一招用了两遍：\textbf{(i)} 造一个逼近目标（无穷大／上确界）的序列；\textbf{(ii)} 用\emph{紧致性}抽出收敛子列，把``逼近''钉在 $K$ 内的一个真实极限点上；\textbf{(iii)} 用\emph{连续性}让函数值跟着收敛过去，于是``逼近的极限''变成``在某点真正达到''。紧致性负责``极限点不逃逸''，连续性负责``函数值不跳变'',二者缺一，达到性就垮。下面两图正是这两根支柱各自抽掉后会怎样。
}

\subsection{两个条件都不能省}

\paragraph{紧致性不可省：开区间上的反例。}
回到引子里的例子。取 $S=(0,1)$（有界但\emph{不闭}，故不紧致），$f(x)=x$ 连续。$f$ 的取值充满 $(0,1)$，上确界 $\sup f=1$，可没有任何 $x\in(0,1)$ 让 $f(x)=1$（图~\ref{fig:compact-2}）。证明里第二步会失败：逼近 $\sup$ 的序列 $x_k\to 1$ 的``极限点'' $1$ 恰恰\emph{漏在 $S$ 之外}——这正是 $S$ 不闭（不紧致）的代价。把定义域换成闭区间 $[0,1]$，端点补回，最大值立刻在 $x^*=1$ 取到。

\begin{figure}[h]
\centering
\begin{tikzpicture}[scale=2.8,>=Stealth]
  % 坐标轴
  \draw[->] (-0.1,0) -- (1.40,0) node[right] {$x$};
  \draw[->] (0,-0.1) -- (0,1.35) node[above] {$f(x)$};
  % 函数 f(x)=x 在开区间 (0,1) 上
  \draw[very thick,blue!60!black] (0,0) -- (1,1);
  % 端点：开圆圈（空心）表示两端不取
  \draw[blue!60!black,fill=white,thick] (0,0) circle (1.6pt);
  \draw[blue!60!black,fill=white,thick] (1,1) circle (1.6pt);
  % 上确界水平虚线 y=1
  \draw[dashed,red!70!black] (0,1) -- (1,1);
  \node[red!70!black,anchor=east] at (-0.04,1) {$1$};
  % x 轴刻度 1
  \node[anchor=north east] at (1,-0.04) {$1$};
  % 标注 sup 取不到
  \node[red!70!black,anchor=west,align=left] at (1.07,1.0) {$\sup f=1$};
  \node[red!70!black,anchor=west,align=left,font=\small] at (1.07,0.84)
     {（无 $x$ 使 $f(x)=1$）};
  \draw[red!70!black,->,thin] (1.05,1.02) -- (1.01,1.0);
  % 函数曲线标签
  \node[blue!60!black,anchor=north west] at (0.40,0.34) {$f(x)=x$};
  % 一串趋近端点但永远到不了的点
  \foreach \px in {0.78,0.88,0.95}{
     \fill[blue!40!black] (\px,\px) circle (0.8pt);
  }
\end{tikzpicture}
\caption{紧致性不可省。$f(x)=x$ 在\emph{开}区间 $(0,1)$ 上连续，取值可任意接近 $\sup f=1$（空心端点处），却没有任何一点使 $f$ 达到 $1$——最大值不存在。}
\label{fig:compact-2}
\end{figure}

\paragraph{连续性不可省。}
反过来，把定义域换成漂亮的紧致集 $[0,1]$，但让 $f$ 不连续：例如
\[
    f(x)=
    \bc
        x, & 0\le x<1,\\[2pt]
        0, & x=1 .
    \ec
\]
$f$ 在 $x=1$ 处跳了一下（图~\ref{fig:compact-3}）。它的上确界仍是 $\sup f=1$，逼近 $\sup$ 的序列 $x_k\to 1$ 极限点 $1$ 这回\emph{留在}定义域内（$K$ 紧致），但证明第二步还是失败：连续性断了，$f(x_k)\to 1$ 与 $f(1)=0$ 对不上，``函数值跟着收敛''这一环掉链子，于是 $f(1)\ne\sup f$，最大值仍不存在。这说明紧致性与连续性各司其职、互不替代。

\begin{figure}[h]
\centering
\begin{tikzpicture}[scale=2.8,>=Stealth]
  % 坐标轴
  \draw[->] (-0.1,0) -- (1.55,0) node[right] {$x$};
  \draw[->] (0,-0.1) -- (0,1.35) node[above] {$f(x)$};
  % 函数 f(x)=x 在 [0,1) 上
  \draw[very thick,blue!60!black] (0,0) -- (1,1);
  % 左端点 (0,0)：实心（含在定义域、连续）
  \fill[blue!60!black] (0,0) circle (1.4pt);
  % x=1 处跳变：上方 (1,1) 空心（f 不取 1），下方 (1,0) 实心（f(1)=0）
  \draw[blue!60!black,fill=white,thick] (1,1) circle (1.6pt);
  \fill[blue!60!black] (1,0) circle (1.6pt);
  % 跳变竖直虚线，连接 (1,0) 与 (1,1)，示意「跳了一下」
  \draw[blue!40!black,dotted,thick] (1,0.06) -- (1,0.94);
  % 上确界水平虚线 y=1
  \draw[dashed,red!70!black] (0,1) -- (1,1);
  \node[red!70!black,anchor=east] at (-0.04,1) {$1$};
  % x 轴刻度 1（放在实心点左下方，避开点与轴箭头）
  \node[anchor=north east] at (0.97,-0.03) {$1$};
  % 标注 sup（右上，避开空心端点与虚线）
  \node[red!70!black,anchor=west,align=left] at (1.10,1.02) {$\sup f=1$};
  \node[red!70!black,anchor=west,align=left,font=\small] at (1.10,0.86)
     {（仍取不到）};
  \draw[red!70!black,->,thin] (1.07,1.0) -- (1.02,1.0);
  % 标注 f(1)=0（放轴下方，远离 x 轴箭头与刻度）
  \node[blue!60!black,anchor=north,align=center,font=\small] at (1.0,-0.13) {$f(1)=0$};
  \draw[blue!60!black,->,thin] (1.0,-0.10) -- (1.0,-0.015);
  % 函数曲线标签
  \node[blue!60!black,anchor=north west] at (0.38,0.36) {$f(x)=x$};
  % 一串趋近 (1,1) 但跳走的点
  \foreach \px in {0.78,0.88,0.95}{
     \fill[blue!40!black] (\px,\px) circle (0.8pt);
  }
\end{tikzpicture}
\caption{连续性不可省。定义域 $[0,1]$ \emph{紧致}，但 $f$ 在 $x=1$ 处跳变：取值可任意接近 $\sup f=1$，而 $f(1)=0$。逼近 $\sup$ 的序列极限点 $1$ 这回留在域内，是\emph{连续性}断了让``函数值跟着收敛''失效——最大值仍不存在。与图~\ref{fig:compact-2} 对照：那里漏的是端点，这里断的是连续。}
\label{fig:compact-3}
\end{figure}

\ex[预算集上的效用最大化为何有解]{
价格 $p\in\RR^n_{++}$（各分量为正）、收入 $w>0$，预算集
\[
    B(p,w)=\setb{x\in\RR^n}{x\ge 0,\ p\cdot x\le w} .
\]
设效用 $u:\RR^n_+\to\RR$ 连续。说明 $\max_{x\in B(p,w)} u(x)$ 存在。
}
\sol{
只需验证 $B(p,w)$ 非空紧致，再调用 Weierstrass 定理。\textbf{非空}：$x=0$ 满足约束。\textbf{有界}：对任一 $x\in B(p,w)$ 与任一坐标 $i$，由 $x\ge 0$ 与 $p\cdot x\le w$ 得 $p_i x_i\le p\cdot x\le w$，故 $0\le x_i\le w/p_i$（这里用到 $p_i>0$）。于是 $B(p,w)$ 被装进盒子 $\prod_i[0,\,w/p_i]$，有界。\textbf{闭}：$B(p,w)$ 是有限个闭集
\[
    \set{x_i\ge 0}\ (i=1,\dots,n)
    \quad\text{与}\quad
    \set{p\cdot x\le w}
\]
的交（每个都是连续函数的非严格不等式水平集，故闭），有限交仍闭。综上 $B(p,w)$ 非空、有界、闭，即非空紧致；$u$ 连续，故由 Weierstrass 定理最大值存在——\emph{马歇尔需求非空、效用最大化问题良定义}。注意这里处处用到 $p_i>0$：若某个价格为零，预算集在该方向无界，紧致性垮掉，``最优消费''可能不存在（想要多少要多少）。
}

\rmkb[只要``达到''，连可微都不需要]{
Weierstrass 定理对 $f$ 只要求\emph{连续}，不要求可微、更不要求凹。这一点常被忽视，却很重要：它把``最优解\emph{存在}''与``用一阶条件\emph{刻画}最优解''彻底分了家。存在性这一关，靠的纯是紧致 $+$ 连续这套拓扑—分析的论证，跟微分毫无关系。等存在性落定，我们才在内点、可微、可能加凹性的额外假设下，用一阶条件去\emph{找}那个早已确知存在的解。
}

\section{一致连续（简述）}
\label{sec:compact-uniform}

连续是``逐点''的性质：在 $x_0$ 处那个 $\delta$ 可以随 $x_0$ 变，越靠近函数``陡''的地方，要达到同样的输出精度 $\ve$，所需的 $\delta$ 往往越小。\textbf{一致连续}则要求一个``通用''的 $\delta$ 对全定义域一齐管用。

\defn{一致连续}{
$f:S\to\RR$ 在 $S$ 上\textbf{一致连续}，若对任意 $\ve>0$，存在\emph{同一个} $\delta>0$，使得对\emph{一切} $x,y\in S$，
\[
    \norm{x-y}<\delta
    \;\To\;
    \abs{f(x)-f(y)}<\ve .
\]
}

与逐点连续的差别全在量词次序：逐点连续是``$\forall x_0\,\forall\ve\,\exists\delta$''（$\delta$ 可依赖 $x_0$），一致连续是``$\forall\ve\,\exists\delta\,\forall x,y$''（$\delta$ 对所有点统一）。$f(x)=1/x$ 在 $(0,1)$ 上逐点连续却\emph{非}一致连续——越靠近 $0$ 越陡，没有一个 $\delta$ 能通吃。但只要把定义域换成紧致集，这种``陡到失控''就不会发生。

\thm{Heine--Cantor 定理}{
若 $K\se\RR^n$ 紧致，$f:K\to\RR$ 连续，则 $f$ 在 $K$ 上\emph{一致}连续。
}

\rmk{证明仍是``紧致 $+$ 连续''那套老配方：反设非一致连续，则可造出两列点 $x_k,y_k$ 距离趋于零而函数值之差始终 $\ge\ve_0$；用序列紧致抽收敛子列把它们逼到同一个极限点，再用连续性导出函数值之差应趋于零——矛盾。}

\state{一致连续以后用在哪}{
它是把``逐点''升级为``整体''的桥梁，在后续两处反复出现：其一，证明连续函数在闭区间上\textbf{可积}（Riemann 积分存在）时，要用一致连续保证分割足够细时各小段的振幅一齐变小（见后文积分与期望部分）；其二，在渐近统计里证明\textbf{一致大数定律}、进而证明极值估计量\emph{一致}时，需要目标函数关于参数``一致''地连续——紧致参数空间正是这条性质的来源。这又是一例：紧致性默默地把``逐点成立''兜底成``整体成立''。
}

\section{这把工具通向何处}
\label{sec:compact-beyond}

本章给出的是经济学里\emph{所有}存在性论证的源头。它的去处密布在三门核心课中：

\begin{itemize}
    \item \textbf{最大值定理}（Berge，见后文紧接的一章）。当目标函数还依赖参数、可行集本身也随参数移动（如预算集随价格变）时，我们不仅要最优\emph{值}存在，还想知道它和最优\emph{选择}如何\emph{连续地}随参数变化。Berge 最大值定理正是 Weierstrass 的参数化升级版，它对约束集（一个对应／集值映射）施加紧致性与（半）连续性，结论是值函数连续、最优解对应上半连续——需求函数的连续性、动态规划里值函数的连续性，都从这里来。
    \item \textbf{均衡的存在性}（不动点定理，见后文）。Brouwer 与 Kakutani 不动点定理的前提里都端坐着``非空、紧致、凸''——紧致性是其中绕不开的一环。一般均衡的存在、纳什均衡的存在，归根结底都建立在``在紧致集上做事''之上。
    \item \textbf{极值估计量的良定义}（见后文渐近统计部分）。极大似然、GMM、M-估计量都形如 $\hat\theta=\argmax_{\theta\in\Theta} Q_n(\theta)$。要让这个 $\argmax$ \emph{存在}（估计量良定义）、并进而\emph{一致}（$\hat\theta\pto\theta_0$），标准做法正是假定参数空间 $\Theta$ 紧致、目标 $Q_n$（及其极限）连续，然后套用本章的 Weierstrass 与一致连续——计量教科书里``$\Theta$ compact''这条看似形式化的假设，兑现的就是本章的存在性。
\end{itemize}

把这几处连起来看，本章其实只讲了一句话：\emph{紧致 $+$ 连续 $\To$ 最优解存在}。可正是这句话，撑起了从消费者选择到一般均衡、再到计量估计的整条``先有解、再谈性质''的链条。先确认解存在，是一切最优化分析当之无愧的第一步。
